\documentclass[11pt,a4paper]{article}

% ==============================================================================
% PREAMBLE: MODERN TECHNICAL ARTICLE STYLE (latex-notes-generator)
% ==============================================================================

% 1. Page Geometry & Typography
\usepackage[margin=2.5cm]{geometry}
\usepackage{mathpazo}
\usepackage{microtype}
\usepackage{setspace}
\setstretch{1.15}

% 2. Math & Symbols
\usepackage{amsmath, amssymb}

% 3. List spacing
\usepackage{enumitem}
\setlist{itemsep=0.4em, topsep=0.4em}

% 4. Callout boxes (tcolorbox)
\usepackage[many]{tcolorbox}
\newtcolorbox{definitionbox}[1][]{
    colback=blue!3!white,
    colframe=blue!60!black,
    fonttitle=\bfseries,
    title=#1,
    arc=2mm
}
\newtcolorbox{formulabox}{
    colback=gray!5!white,
    colframe=gray!40!black,
    boxrule=0.5pt,
    arc=1mm
}

% 5. Header & Footer
\usepackage{fancyhdr}
\pagestyle{fancy}
\fancyhf{}
\fancyhead[L]{\textbf{Lab 01 --- Exploratory Data Analysis}}
\fancyhead[R]{ML \& Data Mining --- UniBo}
\fancyfoot[C]{\thepage}
\renewcommand{\headrulewidth}{0.4pt}

% 6. Hyperlinks
\usepackage{hyperref}
\hypersetup{
    colorlinks=true,
    linkcolor=blue!70!black,
    urlcolor=blue!70!black,
    citecolor=green!50!black
}

% 7. Code listings configuration
\usepackage{xcolor}
\usepackage{listings}
\definecolor{codegreen}{rgb}{0.12, 0.55, 0.15}
\definecolor{codegray}{rgb}{0.45, 0.45, 0.45}
\definecolor{codepurple}{rgb}{0.55, 0.10, 0.70}
\definecolor{codeblue}{rgb}{0.05, 0.25, 0.75}
\definecolor{backcolour}{rgb}{0.97, 0.98, 0.99}
\definecolor{boxframe}{rgb}{0.80, 0.85, 0.90}

\lstdefinestyle{moderncode}{
    backgroundcolor=\color{backcolour},
    commentstyle=\color{codegreen}\itshape,
    keywordstyle=\color{codeblue}\bfseries,
    numberstyle=\tiny\color{codegray},
    stringstyle=\color{codepurple},
    basicstyle=\ttfamily\footnotesize,
    breaklines=true,
    numbers=left,
    numbersep=6pt,
    tabsize=4,
    frame=single,
    rulecolor=\color{boxframe},
    xleftmargin=1.5em,
    framexleftmargin=1.5em
}
\lstset{style=moderncode}

\begin{document}

\section*{Lab 01: Exploratory Data Analysis \& Data Quality}
\addcontentsline{toc}{section}{Lab 01: Exploratory Data Analysis \& Data Quality}

\noindent\textbf{Course:} Machine Learning and Data Mining --- Prof. Claudio Sartori\\
\textbf{Degree:} Master in Artificial Intelligence, Università di Bologna (UniBo)\\
\textbf{Focus Datasets:} Fisher's Iris Dataset ($\mathbb{R}^{150 \times 4}$, supervised 3-class) \& NordicMart Customer Churn ($\mathbb{R}^{3000 \times 8}$, synthetic CRM)

\vspace{0.8em}
\hrule
\vspace{1.2em}

% ==============================================================================
% 1. CORE MACHINE LEARNING CONCEPTS & THEORY
% ==============================================================================
\section{Core Machine Learning Concepts \& Theoretical Foundations}

Exploratory Data Analysis (EDA) is not merely visual inspection; it is the foundational auditing phase that determines data preprocessing pipelines, valid distance metrics, loss functions, and feature engineering decisions.

\subsection{Stevens' Typology of Measurement Scales (1946)}
Every feature in a feature matrix $\mathbf{X} \in \mathbb{R}^{n \times d}$ belongs to a specific scale of measurement that strictly dictates which mathematical operations and algorithms are mathematically valid.

\begin{definitionbox}[Stevens' Measurement Scales in Machine Learning]
\begin{enumerate}[leftmargin=*]
    \item \textbf{Nominal Scale:} Qualitative categories with no intrinsic ordering ($=\,, \neq$). Permitted statistics: Mode, entropy. Operations like arithmetic mean or distance metrics are invalid.
    \begin{itemize}
        \item \textit{Lab Examples:} `class` in Iris; `customer_id` (identifier); `postcode` (categorical code).
        \item \textit{ML Rule:} Identifiers must be removed before feature modeling to prevent data leakage and arbitrary label memorization.
    \end{itemize}
    \item \textbf{Ordinal Scale:} Distinct categories with a strict, transitive monotonic order ($<, \le, =, \ge, >$). Distances between ranks are undefined. Permitted statistics: Median, percentiles, Spearman rank correlation.
    \begin{itemize}
        \item \textit{Lab Example:} `tier` ($\text{Basic} < \text{Silver} < \text{Gold}$).
        \item \textit{ML Rule:} Encoded via ordinal integer mapping ($0 < 1 < 2$) for monotonic models or one-hot encoding if step-distances are non-uniform.
    \end{itemize}
    \item \textbf{Interval Scale:} Quantitative values where differences are meaningful, but there is no absolute/true physical zero. Ratios are meaningless ($20^\circ\text{C}$ is not ``twice as hot'' as $10^\circ\text{C}$). Permitted statistics: Mean, variance, Pearson correlation.
    \item \textbf{Ratio Scale:} Quantitative physical values with an absolute, physically meaningful zero ($0 \implies$ total absence of the attribute). Ratios are strictly defined ($4\,\text{cm}$ is physically twice $2\,\text{cm}$).
    \begin{itemize}
        \item \textit{Lab Examples:} `sepal_length`, `petal_width`, `annual_spend`, `days_since_last_buy`.
        \item \textit{ML Rule:} Euclidean distance ($\ell_2$), Mahalanobis distance, and power/log transformations are mathematically defined.
    \end{itemize}
\end{enumerate}
\end{definitionbox}

\subsection{Distribution Skewness and Statistical Robustness}
Let $\mathcal{X} = \{x_1, x_2, \dots, x_n\}$ be a univariate sample. Choosing between parametric location/dispersion estimators ($\mu, \sigma$) and non-parametric rank estimators ($\text{Med}, \text{IQR}$) depends on distribution symmetry and the \textbf{breakdown point}.

\begin{formulabox}
\textbf{The Breakdown Point ($\varepsilon^*$):} The maximum proportion of arbitrarily corrupted observations a statistic can tolerate before its estimate can be forced to $\pm \infty$.
\begin{equation}
    \varepsilon^*(\text{Sample Mean}) = 0, \quad \varepsilon^*(\text{Sample Variance}) = 0
\end{equation}
\begin{equation}
    \varepsilon^*(\text{Sample Median}) = 0.5, \quad \varepsilon^*(\text{Interquartile Range}) = 0.5
\end{equation}
\end{formulabox}

When a distribution is right-skewed (such as `annual_spend`, where $\mu = \$775.14 \gg \text{Med} = \$574.67$):
\begin{itemize}
    \item Multiplying a single extreme value by 10 increases the standard deviation from $\$673.81$ to $\$1193.07$ ($+77\%$), whereas the median ($\$574.67$) and IQR ($\$594.06$) remain mathematically identical.
    \item \textbf{Implication for ML:} Linear models and gradient descent are vulnerable to leverage points from non-robust features. Skewed ratio variables must be stabilized using monotonic concave transformations such as $\log_{10}(x + 1)$ or Box-Cox.
\end{itemize}

\subsection{Tukey's IQR Outlier Rule and Subgroup Conditioning}
Under Tukey's boxplot rule, candidate outliers are points lying outside the inner fences:
\begin{equation}
    \text{Fences} = \left[ Q_1 - 1.5 \cdot \text{IQR},\; Q_3 + 1.5 \cdot \text{IQR} \right]
\end{equation}

\begin{definitionbox}[Outlier Fallacy: Rare but Real vs. Corrupted Data]
A statistical outlier according to an aggregated IQR fence is \textbf{not necessarily an error}:
\begin{itemize}
    \item In NordicMart, 204 customers exceed the upper fence ($\$1843.24$). However, stratifying by `tier` reveals that 131 belong to the Gold tier ($64\%$). For Gold customers, $Q_3 = \$1909.85$ and the median is $\$1297.18$.
    \item \textbf{Takeaway:} These data points are genuine high-value customers (``whales''), not measurement errors. Truncating or deleting them removes the most profitable segment from the training data.
\end{itemize}
\end{definitionbox}

\subsection{Simpson's Paradox and Confounder Variables}
Simpson's Paradox occurs when an apparent statistical association observed across aggregated data reverses or disappears when the data is conditioned on a confounding variable $Z$.

\begin{formulabox}
Mathematically, for random variables $X, Y$ and confounder $Z$:
\begin{equation}
    \text{Cov}(X, Y) \neq \sum_{z} P(Z = z) \cdot \text{Cov}(X, Y \mid Z = z)
\end{equation}
\end{formulabox}

\begin{itemize}
    \item \textbf{Iris Example:} Globally, $\text{Corr}(\text{sepal\_length}, \text{sepal\_width}) = -0.12$ (slightly negative). Conditioned within species:
    \begin{equation}
        r_{\text{setosa}} = +0.74, \quad r_{\text{versicolor}} = +0.53, \quad r_{\text{virginica}} = +0.46
    \end{equation}
    The species is a confounder: *setosa* flowers have short but very wide sepals, whereas the other species have longer but narrower sepals.
    \item \textbf{NordicMart Example:} Globally, $\text{Corr}(\text{annual\_spend}, \text{web\_sessions\_month}) = +0.37$. Within tiers:
    \begin{equation}
        r_{\text{Basic}} = +0.04, \quad r_{\text{Silver}} = +0.04, \quad r_{\text{Gold}} = -0.00
    \end{equation}
    Browsing the website does \textbf{not} cause higher spending. Loyalty tier is the common cause: Gold customers independently spend more and visit the website more frequently.
\end{itemize}

\subsection{Taxonomy of Missing Data (Rubin, 1976)}
Handling missing values ($\text{NaN}$) depends critically on the missingness mechanism:
\begin{enumerate}[leftmargin=*]
    \item \textbf{MCAR (Missing Completely at Random):} $P(\text{Missing} \mid X, Y) = P(\text{Missing})$. Missingness is unrelated to observed or unobserved data.
    \begin{itemize}
        \item \textit{Diagnosis in Lab:} For `age`, churn rate with missing age is $12.6\%$ vs $11.9\%$ when observed (statistically indistinguishable). Imputation by median is safe and unbiased.
    \end{itemize}
    \item \textbf{MAR (Missing at Random):} $P(\text{Missing} \mid X, Y) = P(\text{Missing} \mid X)$. Missingness depends on observed features but not on the missing value itself.
    \item \textbf{MNAR (Missing Not at Random / Informative Missingness):} $P(\text{Missing} \mid X, Y)$ depends directly on the unobserved variable itself or target $Y$.
    \begin{itemize}
        \item \textit{Diagnosis in Lab:} For `web_sessions_month`, customers with missing values have a churn rate of \textbf{$23.1\%$}, compared to \textbf{$9.6\%$} for observed sessions ($2.4\times$ higher risk).
        \item \textit{Action:} Missingness is an informative behavioral feature. We must create an explicit binary indicator flag `web_sessions_isna` prior to imputing values.
    \end{itemize}
\end{enumerate}

\subsection{Class Imbalance and the Fallacy of Raw Accuracy}
In binary classification with target $y \in \{0, 1\}$, where the positive class is rare:
\begin{formulabox}
\textbf{The Trivial Majority-Class Classifier:}
\begin{equation}
    \hat{y}_{\text{trivial}} = \arg\max_{c \in \{0, 1\}} P(y = c) = 0 \quad (\forall i)
\end{equation}
In NordicMart, non-churners represent $88.0\%$ ($n_0 = 2640$) and churners represent $12.0\%$ ($n_1 = 360$).
\begin{equation}
    \text{Accuracy}(\hat{y}_{\text{trivial}}) = 88.0\%, \quad \text{Recall}(\hat{y}_{\text{trivial}}) = \frac{0}{360} = 0.0\%
\end{equation}
\end{formulabox}
An accuracy of $88\%$ provides zero commercial value. Models trained on imbalanced data must be evaluated using Precision, Recall, $F_1$-score, PR-AUC, or cost-utility matrices.

\vspace{1em}
\newpage

% ==============================================================================
% 2. CANONICAL CODE SNIPPETS & EDA RECIPES
% ==============================================================================
\section{Canonical Code Snippets \& EDA Recipes}

\subsection{Recipe 1: Data Audit, Identifier Isolation \& Ordered Categoricals}
\begin{lstlisting}[language=Python]
import pandas as pd
import numpy as np

# 1. Load raw data (never mutate df directly)
df = pd.read_csv("nordicmart_customers.csv")
customers = df.copy()

# 2. Audit cardinality: detect constants and unique identifiers
cardinality = customers.nunique()
identifiers = cardinality[cardinality == len(customers)].index.tolist()
constants = cardinality[cardinality <= 1].index.tolist()

# 3. Clean string codes and enforce ordinal hierarchies
customers["postcode"] = customers["postcode"].astype("string").str.replace(r"\.0$", "", regex=True)
customers["tier"] = pd.Categorical(
    customers["tier"],
    categories=["Basic", "Silver", "Gold"],
    ordered=True
)
\end{lstlisting}

\subsection{Recipe 2: Sentinel Value Extraction \& Binary Indicator Generation}
\begin{lstlisting}[language=Python]
# Detect hidden database sentinels (e.g. 999 days, 0 age)
# Step 1: Create explicit indicator feature before recoding
customers["never_bought"] = (customers["days_since_last_buy"] == 999).astype(int)

# Step 2: Replace sentinel values with NaN (avoid deprecated inplace=True)
customers["days_since_last_buy"] = customers["days_since_last_buy"].replace(999, np.nan)
customers["age"] = customers["age"].replace(0, np.nan)
\end{lstlisting}

\subsection{Recipe 3: Informative Missingness Diagnostic}
\begin{lstlisting}[language=Python]
# Diagnostic: Check if missingness predicts the target
missing_diagnostic = customers.groupby(customers["web_sessions_month"].isna())["churned"].agg(
    churn_rate="mean",
    count="count"
)
print(missing_diagnostic)
# If churn rate differs significantly between True and False -> MNAR -> Keep indicator flag!
\end{lstlisting}

\subsection{Recipe 4: Subgroup-Stratified Correlation (Checking Confounders)}
\begin{lstlisting}[language=Python]
# Check global Pearson correlation vs intra-group correlation
r_global = customers["annual_spend"].corr(customers["web_sessions_month"])
print(f"Global correlation: {r_global:.2f}")

print("Intra-group correlations by tier:")
for tier_name, group in customers.groupby("tier", observed=False):
    r_sub = group["annual_spend"].corr(group["web_sessions_month"])
    print(f"  Tier {tier_name:6s}: {r_sub:.2f}")
\end{lstlisting}

\subsection{Recipe 5: Bivariate Evaluation of Numeric Features vs Target}
\begin{lstlisting}[language=Python]
import matplotlib.pyplot as plt
import seaborn as sns

fig, axs = plt.subplots(1, 2, figsize=(12, 4))
# Comparing discriminative capacity across classes
sns.boxplot(data=customers, x="churned", y="days_since_last_buy", ax=axs[0])
axs[0].set_title("Recency by Churn (High separation)")

sns.boxplot(data=customers, x="churned", y="tickets_90d", ax=axs[1])
axs[1].set_title("Tickets by Churn (Negligible separation)")
plt.tight_layout()
plt.show()
\end{lstlisting}

\vspace{1em}

% ==============================================================================
% 3. PITFALLS, TRAPS & ANTI-PATTERNS
% ==============================================================================
\section{Pitfalls, Traps \& Anti-Patterns}

\begin{enumerate}[leftmargin=*]
    \item \textbf{Trapping Identifiers in Feature Sets (Data Leakage):}\\
    \textit{Anti-pattern:} Passing \lstinline!customer_id! into Pearson correlation matrices or decision trees.\\
    \textit{Consequence:} In a decision tree, an identifier with cardinality equal to $N$ can split the entire dataset perfectly, producing training accuracy $100\%$ but zero generalization capability.
    
    \item \textbf{Confusing Column Duplicates with Row Duplicates:}\\
    \textit{Anti-pattern:} Calling \lstinline!customers["col"].duplicated().sum()! instead of \lstinline!customers.duplicated().sum()!.\\
    \textit{Consequence:} Repeating values within a single continuous or categorical column is standard. A duplicated record exists only when \textbf{all features simultaneously match} across two rows.

    \item \textbf{The In-Place Mutation Trap (`inplace=True`):}\\
    \textit{Anti-pattern:} \lstinline!customers["age"].replace(0, np.nan, inplace=True)!, which is deprecated in modern pandas (raises \texttt{FutureWarning} or \texttt{SettingWithCopyWarning}).\\
    \textit{Best practice:} Always use explicit functional reassignment: \lstinline!customers["age"] = customers["age"].replace(0, np.nan)!.

    \item \textbf{Naive Imputation under Informative Missingness (MNAR):}\\
    \textit{Anti-pattern:} Directly replacing missing values in \lstinline!web_sessions_month! with the column median without generating a missing indicator flag.\\
    \textit{Consequence:} Eliminates the strongest predictive signal in the dataset ($23.1\%$ churn vs $9.6\%$).
\end{enumerate}

\vspace{1em}

% ==============================================================================
% 4. MASTER-LEVEL EXAM & DEFENSE QUESTIONS
% ==============================================================================
\section{Master-Level Exam \& Defense Questions}

\begin{definitionbox}[Question 1: Mathematical Mechanism of Simpson's Paradox]
\textbf{Question:} Formally explain how an overall correlation between two continuous variables can be positive while the correlation within every subgroup is zero or negative. How should an ML engineer detect this during EDA before training a regression model?\\
\textbf{Answer:} By the law of total covariance:
\[
\text{Cov}(X, Y) = \mathbb{E}_Z[\text{Cov}(X, Y \mid Z)] + \text{Cov}_Z(\mathbb{E}[X \mid Z], \mathbb{E}[Y \mid Z])
\]
Even if within-group covariance $\text{Cov}(X, Y \mid Z = z) \approx 0$ for all $z$, the between-group covariance $\text{Cov}_Z(\mathbb{E}[X \mid Z], \mathbb{E}[Y \mid Z])$ can be large and positive if group centroids move jointly in the $(X, Y)$ plane. In NordicMart, $\mathbb{E}[\text{spend} \mid \text{Gold}] > \mathbb{E}[\text{spend} \mid \text{Basic}]$ and $\mathbb{E}[\text{sessions} \mid \text{Gold}] > \mathbb{E}[\text{sessions} \mid \text{Basic}]$. An ML engineer detects this by generating conditional scatter plots colored by categorical features and computing intra-group correlation matrices across all nominal/ordinal variables.
\end{definitionbox}

\begin{definitionbox}[Question 2: Pearson vs. Spearman Divergence on Skewed Targets]
\textbf{Question:} In NordicMart, the correlation between \lstinline!days_since_last_buy! and \lstinline!churned! is $r = 0.38$ under Pearson, but drops to $\rho = 0.28$ under Spearman. Why does this divergence occur, and which metric provides a more faithful assessment of predictive utility?\\
\textbf{Answer:} Pearson correlation assesses the degree of strictly linear association and is heavily weighted by points in the extreme right tail of \lstinline!days_since_last_buy! due to squared deviation terms $(x_i - \bar{x})^2$. Spearman correlation replaces raw numerical values with their fractional ranks ($1$ to $n$), eliminating the disproportionate leverage of extreme recency values. Because machine learning models (especially decision trees and non-linear kernels) depend on monotonic ordering rather than strict linearity, Spearman $\rho$ gives a more robust estimate of whether sorting instances by recency separates churners.
\end{definitionbox}

\begin{definitionbox}[Question 3: Handling Informative Missingness in Production Pipelines]
\textbf{Question:} Suppose you must deploy a logistic regression model on NordicMart data. How should you treat \lstinline!web_sessions_month! given that missingness is MNAR and represents $18\%$ of rows?\\
\textbf{Answer:} 
1. \textbf{Missingness Indicator Feature:} Create a binary dummy feature $M_{\text{web}} = \mathbf{1}_{\{\text{web\_sessions\_month} = \text{NaN}\}}$.
2. \textbf{Imputation:} Impute the missing values in the continuous feature with a central location estimate (e.g. median of observed values) or a constant (e.g. $0$).
3. \textbf{Regression Formulation:} The linear model becomes:
\[
\hat{y} = \sigma\left(\beta_0 + \beta_1 x_{\text{web}} + \beta_2 M_{\text{web}} + \dots\right)
\]
When $M_{\text{web}} = 1$, $\beta_2$ explicitly captures the baseline shift in log-odds associated with missingness, preventing the imputed constant from corrupting the estimation of $\beta_1$.
\end{definitionbox}

\begin{definitionbox}[Question 4: Proof of Zero Breakdown Point for the Sample Mean]
\textbf{Question:} Prove mathematically that the breakdown point of the sample mean $\bar{x} = \frac{1}{n} \sum_{i=1}^n x_i$ is $\varepsilon^* = 0$, whereas the sample median has $\varepsilon^* = 0.5$.\\
\textbf{Answer:} 
Let $\mathcal{X} = \{x_1, \dots, x_n\}$ and replace a single observation $x_1$ with an arbitrary corrupted value $x_1^* \in \mathbb{R}$. The corrupted mean is:
\[
\bar{x}^* = \frac{x_1^* + \sum_{i=2}^n x_i}{n} \implies \lim_{x_1^* \to \infty} |\bar{x} - \bar{x}^*| = \infty
\]
Since a single corrupted data point ($m = 1$) out of $n$ can drive the estimate to infinity:
\[
\varepsilon^* = \lim_{n \to \infty} \frac{1}{n} = 0
\]
For the sample median, if $m < \lfloor \frac{n+1}{2} \rfloor$ points are corrupted, the median is constrained by the remaining uncorrupted order statistics $x_{(k)}$. Thus, $\varepsilon^* = \frac{\lfloor (n-1)/2 \rfloor}{n} \to 0.5$ as $n \to \infty$.
\end{definitionbox}

\end{document}
